\section*{The Arthur--Merlin protocol (author Section 5)}
The starting point for this section is the conclusion of Theorem 18:
\paragraph{Hypothesis AM$(N)$.} $N$ is an oracle machine such such that for all $\gamma\in[0,1)$, $N^{J[\gamma]}(x)$ computes $L$ with advantage $\epsilon/2$.
We will show that this implies that there is an Arthur--Merlin protocol for $L$ whose running time is in the running time of $N$.
We make use of two known AM protocols. The input to both protocols includes a function $f:\{0,1\}^w\longrightarrow\{0,1\}^*$ given by a circuit $C_f$. The function $f$ induces a probability distribution $\pi_f$ on $\{0,1\}^*$, with $\pi_f(q)$ equal to $|f^{-1}(q)|/2^2$, which is the probability that a uniformly chosen element of $\{0,1\}^w$ maps to $q$. Also $\zeta,\delta>0$ are input parameters for the protocol.
\paragraph{Lemma 21 (Entropy Estimation Protocol, [24]).} There is a constant round AM-protocol that takes as input a circuit $C_f$ as above, an integer $h$ and an error parameter $\delta$, runs in time $|C_f|\operatorname{poly}(\log(1/\delta))$, and outputs accept or reject with the following conditions:
\begin{enumerate}\item If $|h-H(\pi_f)|\le1$ then Merlin has a strategy that causes Arthur to reject with probability at most $\delta$.
\item If $|h-H(\pi_f)|\ge2$ then the probabilty that the verification is accepted is at most $\delta$.\end{enumerate}
\paragraph{Lemma 22 (Lower Bound protocol, [25]).} There is a constant round protocol with input $C_f$, input sequences $q_1,\ldots,q_m\in(\{0,1\}^*)^m$ and $(s_1,\ldots,s_m)\in\mathbb N^m$, and error parameters $(\zeta,\delta)$ both in $(0,1)$ that runs in time $C_f\operatorname{poly}(m,1/\zeta,\log(1/\delta))$, and outputs either accept or reject according to the following:
\begin{enumerate}\item If $s_i\ge|f^{-1}(q_i)|$ for every $i$, then Merlin has a strategy such that the protocol rejects with probability at most $\delta$.
\item If $s_i\le|f^{-1}(q_i)|/(1+\zeta)$ for some $i$ then for any strategy of Merlin the probabilty that the verification is accepted is at most $\delta$.\end{enumerate}
We fix some notation for this section: $x$ denotes the string whose membership in $L$ is to be determined. $r=r_x^N$ is the length of the random string generated by $N$ on input $x$. $k=k_x^N$ is the number of queries asked by $N$ on input $x$. We assume (with no significant loss of generality) that $\log k$ is an integer. For $\rho\in\{0,1\}^r$, $q_1(\rho),\ldots,q_k(\rho)$ denotes the queries asked by $N$ on input $x$. $f$ denotes the function on $\{0,1\}^{r+\log k}$ that maps $(\rho,i)$ to $q_i(\rho)$. Let $\pi=\pi_f$ be the probability distribution on $\{0,1\}^*$ induced by $f$, which is given by $\pi_f=|f^{-1}(q)|/2^{r+\log k}$. $Q=Q_x^N$ is the image of the function $f$, which is the set of queries that are asked with positive probability. For $\rho\in\{0,1\}^r$, $a_1[\gamma](\rho),\ldots,a_k[\gamma](\rho)$ is the sequence of answers from the oracle $J[\gamma]$ to the queries $q_1(\rho),\ldots,q_k(\rho)$. When the random string is $\rho$, the output of the algorithm is $\operatorname{out}(a_1[\gamma](\rho),\ldots,a_k[\gamma](\rho);\rho)$. We denote this by $z[\gamma](\rho)$. $\epsilon=\epsilon_{|x|}^N$.
\paragraph{Proposition 23.} Let $\pi$ be a probability distribution over $\{0,1\}^w$.
\begin{enumerate}\item If $q$ is selected according to $\pi$, $\mathbb E[\log(1/\pi(q))]=H(\pi)$ where $H(\pi)$ is the binary entropy of $\pi$.
\item Suppose $q_1,\ldots,q_t$ are selected independently according to $\pi$. Then
\[\Pr\left[\left|\frac1t\sum_{i=1}^t\log(1/\pi(q_i))-H(\pi)\right|>\zeta\right]\le\frac{w^2}{t\zeta^2}.\]\end{enumerate}
\begin{proof}The first part follows directly from the definition of $H(\pi)$. For the second part, let $Z_i=\log(1/\pi(q_i))$ and let $Z=\frac1t\sum Z_i$. Note that $\mathbb E[Z]=\frac1t\sum\mathbb E[Z_i]=H(\pi)$ by the first part. Also, note that since $Z$ is an average of $t$ i.i.d. random variables, it's variance is $\frac1t\operatorname{Var}(Z_i)\le w^2/t$ since a random variable taking values in $[0,w]$ has variance at most $w^2$. By Chebyshev's inequality, for any random variable $X$, the probability that $|Z-\mathbb E[Z]|>\zeta$ is at most $\operatorname{Var}(Z)/\zeta^2\le w^2/(t\zeta^2)$.\end{proof}
In describing the Arthur--Merlin protocol we refer to the actions of ``Honest Merlin'' which is the ideal behavior for Merlin (which Merlin may not follow). The protocol makes use of several parameters, whose values are determined by the analysis. For reference we list these parameters here with their values:
\begin{itemize}\item $t$ is the sample complexity, and it set to $t=32000\max(r,1)^2k^2/\epsilon(|x|)^2$.
\item $\tau$ is the safety parameter, which is set to $\tau=\epsilon(|x|)/(16(4k+2))$.
\item $\nu$ is the lower bound parameter, which is set to $\nu=\epsilon(|x|)^2/(20000k^2)$.
\item $D$ is the precision parameter and is a positive integer. We choose $D=\lceil16(4k+2)/\epsilon(|x|)\rceil$.\end{itemize}
It is straightforward to verify that these parameter values satisfy the following conditions:
\paragraph{Proposition 24.}
\begin{align*}
\tau&\le\frac{\epsilon(|x|)}{16(4k+2)},& \nu&\le\tau\epsilon(|x|)/(128k),\\
t&\ge\frac{256k}{\epsilon(|x|)}\log\frac{16}{\epsilon(|x|)},& t&\ge\frac{16(r+\log k)^2}{\epsilon(|x|)},\\
t&\ge\frac{320k}{\epsilon(|x|)\tau},&D&\ge\frac1\tau,\\
D&\ge\frac{32(k+1)}{\epsilon(|x|)}.&&
\end{align*}
The inequalities stated in this Proposition are needed in the analysis of the protocol. The values of the parameters were chosen to ensure these conditions. We are now ready to state the protocol:
\paragraph{Step 1.} Merlin provides an integer $h$. Arthur and Merlin perform the Entropy estimation protocol of Lemma 21 for the probability distribution $\pi$ and input $h$ with error parameter $\delta=\epsilon(|x|)/16$. (If Merlin is honest, then Merlin sets $h$ so that $|h-H(\pi)|\le1$, and adopts the strategy from the first part of Lemma 21 that makes the rejection probability at most $\epsilon(|x|)/16$.)
\paragraph{Step 2.} Random strings $\rho_1,\ldots,\rho_t$, each of length $r$, are generated using public coins. Let $\chi$ be the $t\times k$ matrix with $\chi_{i,j}=f(\rho_i,j)$, the $j$th query asked when the random string is $\rho_i$.
\paragraph{Step 3.} Merlin provides a $t$ by $k$ positive integer matrix $B$ (Honest Merlin sets $B$ to be equal to the matrix $B^*$ whose $i,j$ entry is $|f^{-1}(\chi_{i,j})|$.)
\paragraph{Notational Remark.} It is convenient here to define some matrices and vectors that are determined from $B$. For $\gamma\in[0,1)$ the $t\times k$ matrix $A[\gamma]$ is given $A[\gamma]_{i,j}=\phi((1+\gamma)B_{i,j})$, and the vector $z[\gamma]\in\{0,1\}^t$ is given by $z[\gamma]_i=\operatorname{out}(A[\gamma]_{i,1},\ldots,A[\gamma]_{i,k})$. The matrix $A^*[\gamma]$ and vector $z^*[\gamma]$ are obtained analogously with $B$ replaced by $B^*$. Note that row $i$ of $A^*[\gamma]$ is exactly the sequence of query answers provided by oracle $J[\gamma]$ when the random string is $\rho_i$ and $z^*[\gamma]_i$ is the output of $N^{J[\gamma]}(x)$ for random string is $\rho_i$.
\paragraph{Step 4.} For each row index $i$, Arthur randomly chooses $\hat j(i)\in\{1,\ldots,k\}$. Define
\[\sigma(B)=\frac1t\sum_i(r-\log(B_{i,\hat j(i)})).\]
Arthur computes $\sigma(B)$ approximately to determine an integer $H'$ that belongs to the interval $[\sigma(B)-1,\sigma(B)+1]$. Arthur rejects this step if $|H'-h|>3$ and accepts this step otherwise.
\paragraph{Step 5.} Arthur--Merlin perform the lower bound protocol of Lemma 22 for the above-mentioned function $f$ that maps $(\rho,i)$ to $q_i(\rho)$, with error parameters $(\zeta,\delta)=(\nu,\epsilon(|x|)/16)$. The input $(q_1,\ldots,q_m),(s_1,\ldots,s_m)$ to the protocol is the sequence of $tk$ entries of the matrix $\chi$ and the corresponding sequence of entries of the matrix $B$. (If Merlin is Honest then since $B=B^*$, $s_i=|f^{-1}(q_i)|$ for every $i$ and Merlin follows the strategy from the first part of Lemma 22 so that the probability of rejection is at most $\epsilon(|x|)/16$.)
\paragraph{Step 6.} Arthur selects a uniformly random $\hat\lambda\in\{1,\ldots,D\}$ and a uniformly random $\hat i\in\{1,\ldots,t\}$. $\hat z$ is defined to be $z[\lambda/D]_{\hat i}$. For future reference we also define $\hat z^*$ to be $\hat z$ if $B$ is replaced by $B^*$.
\paragraph{Step 7.} If any of the steps 1,4 or 5 result in reject, then the protocol outputs failure. Otherwise the output is $\hat z$.
This protocol outputs either failure or 0 or 1. The main theorem of this subsection is:
\paragraph{Theorem 25.} Under the hypotheses AM$(N)$ given earlier, the above protocol satisfies:
\begin{enumerate}\item If Merlin behaves honestly, then the protocol outputs $L(x)$ with probability at least $\frac12+\epsilon(|x|)/16$.
\item For any behavior of Merlin, the probability that the protocol outputs failure or outputs $L(x)$ is at least $\frac12+\epsilon(|x|)/16$.\end{enumerate}
\paragraph{Corollary 26.} Both $L$ and $\overline L$ can be recognized by AM protocols that run in time polynomial in the $1/\epsilon$ and the running time of $N$.
\begin{proof}The protocol obtained by changing a failure output to 0 (resp. 1) is an AM protocol (resp. coAM protocol) for $L$.\end{proof}
\paragraph{Proof of Theorem 25.} We first note the following:
\paragraph{Proposition 27.} The value $\hat z^*$ determined in Step 6 in the case that $B=B^*$ agrees with $L(x)$ with probability at least $\frac12+\epsilon(|x|)/4$.
\begin{proof}The value $\hat z^*$ determined in Step 6 is equal to the output of $N^{J[\lambda/D]}$ on input $x$ when the random string is $\rho_{\hat i}$. Since $\rho_{\hat i}$ is a uniformly random string, the hypothesis on $N$ implies that this is equal to $L(x)$ with probability at least $\frac12+\epsilon(|x|)/4$.\end{proof}
We now prove the first part of the Theorem. For $s\in\{1,4,5\}$ define the event $R_s$ to be the event that the test in Step $s$ outputs reject. Assume Merlin behaves honestly, and therefore $B=B^*$. By the above proposition, the probability that the output of the protocol is $L(x)$ is at least:
\begin{align*}
&\frac12+\frac{\epsilon(|x|)}4-\Pr[R_1\lor R_4\lor R_5]\\
&\qquad\ge\frac12+\frac{\epsilon(|x|)}4-\Pr[R_1]-\Pr[R_4]-\Pr[R_5].
\end{align*}
As noted in the description of Steps 1 and 5, $\Pr[R_1]$ and $\Pr[R_5]$ are each at most $\epsilon(|x|)/16$. So it suffices to show that $\Pr[R_4]\le\epsilon(|x|)/16$, which we now do. Step 4 is rejected only if $|H'-h|>3$. We have:
\[|H'-h|\le|H'-\sigma(B^*)|+|\sigma(B^*)-H(\pi)|+|H(\pi)-h|.\]
$|H'-\sigma(B^*)|\le1$ by the choice of $H'$ and the fact that $B=B^*$, and $|H(\pi)-h|\le1$ since Merlin is assumed to be honest. Therefore $\Pr[R_4]\le\Pr[|\sigma(B^*)-H(\pi)|>1$. So it suffices to prove:
\paragraph{Proposition 28.} $\Pr[|\sigma(B^*)-H(\pi)|>1]\le\epsilon(|x|)/16$.
\begin{proof}By the definition of $B^*$, $B^*_{i,\hat j}=2^{r+\log k}\pi(\chi_{i,\hat j(i)})$ and so $r+\log k-\log(B^*_{i,\hat j(i)})=\log(1/\pi(\chi_{i,\hat j(i)}))$ and so $\sigma(B^*)=\frac1t\sum_{i=1}^t\log(1/\pi(\chi_{i,\hat j(i)}))$. Since $\chi_{1,\hat j(1)},\ldots,\chi_{t,\hat j(t)}$ is a sequence of independent samples from the distribution $\pi$, Proposition 23 implies
\[\Pr\left[\left|\frac1t\sum_i\log(1/\pi(\chi_{i,\hat j(i)}))-H(\pi)\right|>1\right]<\frac{(r+\log k)^2}t\le\epsilon(|x|)/16,\]
where the last inequality used Proposition 24.\end{proof}
This completes the proof of the first part of the theorem. We now turn to the proof of the second part. We consider an arbitrary strategy for Merlin. We break up into cases according to properties of the matrix $B$. We divide into three cases.
\paragraph{Case 1.} There is at least one entry $(i,j)$ with $B_{i,j}<B^*_{i,j}/(1+\nu)$. By the property of the Lower bound protocol performed in Step 5, this will reject with probability at least $1-\epsilon(|x|)/16\ge\frac12)+\epsilon(|x|)/4$.
\paragraph{Case 2.} $B_{i,j}\ge B^*_{i,j}/(1+\nu)$ for all entries $(i,j)$ and there are more than $\epsilon(|x|)t/16$ entries with $B_{i,j}>(1+\tau)B^*_{i,j}$. We claim that Step 4 will reject with probability at least $1-\epsilon(|x|)/8\ge1/2+\epsilon/4$. For this we give a lower bound on the probability $H'-H(\pi)$ exceeds 3. We have
\begin{align*}H'-H(\pi)&=\sigma(B)-\sigma(B^*)-(\Sigma(B)-H')-(H(\pi)-\sigma(B^*))\\
&\ge\sigma(B)-\sigma(B^*)-|\sigma(B)-H'|-|H(\pi)-\sigma(B^*)|.\end{align*}
By the choice of $H'$, $|\sigma(B)-H'|\le1$. Therefore:
\begin{align*}\Pr[H'-H(\pi)>3]&\ge\Pr[(\sigma(B)-\sigma(B^*)\ge5)\land(|\sigma(B^*)-H(\pi)|<1)]\\
&\ge1-\Pr[\sigma(B)<\sigma(B^*)+5]\\&\qquad-\Pr[|\sigma(B^*)-H(\pi)|\ge1].\end{align*}
The desired lower bound will follow by showing that each of the two probabilities on the right is at most $\epsilon(|x|)/16$. Proposition 28 shows this for the second probability. So we now bound the first probability. For $i\in\{1,\ldots,t\}$ define:
\begin{itemize}\item $V_i=\{j\in\{1,\ldots,k\}:\log B_{i,j}\ge(1+\tau)B^*_{i,j}$. (Note that $\sum_i V_i\ge\epsilon(|x|)t/16$ by the case assumption.)
\item $Y_i$ is the 0-1 indicator of the event that $\hat j(i)\in V_i$, and $Y=\sum_iY_i$.\end{itemize}
We now make two claims. Claim 1. $\sigma(B)-\sigma(B^*)\ge\tau Y-\nu t$. Claim 2. $\Pr[Y<\frac12\frac{\epsilon(|x|)t}{16k}]\le\epsilon(|x|)/16$.
If $Y\ge\frac12\frac{\epsilon(|x|)t}{16k}$, then by the conditions on $\tau$ and $\nu$ given by Proposition 24 we have that $\tau Y-\nu t\ge5$ and so
\[\Pr[\sigma(B)-\sigma(B^*)<5]\le\Pr\left[Y<\frac12\frac{\epsilon(|x|)t}{16k}\right]\le\epsilon(|x|)/16.\]
To prove the first claim, note that for $i$ such that $Y_i=1$, we have
\[\log(B_{i,\hat j(i)})\ge\log((1+\tau)B^*_{i,\hat j(i)}\ge\log(B^*_{i,\hat j(i)})+\tau\]
(since $\log(1+\tau)\ge\tau$ for $\tau\in[0,1]$). For $i$ such that $Y_i=0$, we have $B_{i,\hat j(i)}\ge B_{i,\hat j(i)}/(1+\nu)$ which implies $\log B_{i,\hat j(i)}\ge\log B_{i,\hat j(i)}-2\nu$ since $\log(1+\nu)\le2\nu$). Therefore $\sigma(B)-\sigma(B^*)\ge\tau Y-\nu t$.
For the second claim, note that that $\mathbb E[Y]=\sum_i\mathbb E[Y_i]=\sum_i|V_i|/k\ge\epsilon(|x|)t/(16k)$, by the case assumption. The ``multiplicative Chernoff bound'' implies that if $Y$ is a sum of independent 0-1 random variables, then $\Pr[Y\le\frac12\mathbb E[Y]]\le e^{-\mathbb E[Y]/8}$. Applying this here we obtain $\Pr[Y\le\epsilon(|x|)t/(32k)\le e^{-\epsilon(|x|)t/(256k)}\le\epsilon(|x|)/16$ by Proposition 24.
\paragraph{Case 3.} $B_{i,j}\ge B^*_{i,j}/(1+\nu)$ for all entries $(i,j)$ and there are at most $\epsilon(|x|)t/16$ entries with $B_{i,j}>(1+\tau)B^*_{i,j}$. By Proposition 27, $\hat z^*$ is equal to $L(x)$ with probability at least $\frac12+\epsilon(|x|)/4$. We show that under the case assumption the probability that $\hat z\ne\hat z^*$ is at most $\epsilon(|x|)/8$ which will show that $\hat z$ is equal to $L(x)$ with probability at least $\frac12+\epsilon(|x|)/8$.
We say that a row index $i\in\{1,\ldots,t\}$ is unsafe if there is at least one $j$ such that $B_{i,j}>(1+\tau)B^*_{i,j}$ and is safe otherwise. We identify the following ``bad'' events:
\begin{itemize}\item[B1] The selected row index $\hat i$ is unsafe.
\item[B2] The selected $\hat\lambda$ satisfies
\[\hat\lambda/D\in[0,2\tau)\cup[1-2\tau,1)\cup\bigcup_{j=1}^k[\mu(B^*_{\hat i,j})-2\tau,\mu(B^*_{\hat i,j})+2\tau].\]\end{itemize}
We claim that if neither bad event happens then $\hat z=\hat z^*$. Since [B1] doesn't happen, $\hat i$ is safe. Together with the case assumption and the fact that $\nu\le\tau$ this implies that for each $j\in[k]$, $B_{\hat i,j}\in[B^*_{\hat i,j}/(1+\tau),B^*_{\hat i,j}(1+\tau)]$. Since [B2] doesn't happen, then for all $j\in\{1,\ldots,k\}$, the hypotheses of Proposition 16 hold for $(s,t)=(B^*_{\hat i,j},B_{\hat i,j})$ from which we conclude that for all $j\in\{1,\ldots,k\}$
\begin{align*}A[\lambda/D]_{\hat i,j}&=\phi((1+\lambda/D)B_{\hat i,j}))\\&=\phi((1+\lambda/D)B^*_{\hat i,j}))=A^*[\lambda/D]_{\hat i,j},\end{align*}
which implies that $\hat z[\lambda/D]=\hat z^*[\lambda/D]$. So now it will be enough to show that $\Pr[\mathrm{B1}]$ and $\Pr[\mathrm{B2}]$ are both at most $\epsilon(|x|)/16$. $\Pr[\mathrm{B1}]$ is equal to the fraction of unsafe rows, which is at most $\epsilon(|x|)/16$ by the case assumption.
Since $\lambda$ is chosen uniformly from $\{1,\ldots,D\}$, $\Pr[\mathrm{B2}]$ is the fraction of integers in $\{1,\ldots,D\}$ that belong to the set
\begin{multline*}[0,2D\tau)\cup[D(1-2\tau),D)\\\cup\bigcup_{j=1}^k[D(\mu(B^*_{\hat i,j})-2\tau),D(\mu(B^*_{\hat i,j})+2\tau)].\end{multline*}
Now we use the fact that the number of integers in an interval $[a,b)$ is less than $b-a+1$ and an interval $[a,b]$ is at most $b-a+1$. Therefore the number of integers in $[0,2D\tau)\cup[D(1-2\tau),D)$ is less than $4\tau D+2$ and the number of integers in each set of the big union is at most $4\tau D+1$. Summing up over all intervals the number of integers in the union is at most $(4k+4)\tau D+k+1$. Since $\tau D\le1/4$ by Proposition 24, this is at most $2k+2$. So the probability that $\lambda$ is in this set is at most $(2k+2)/D$ which is at most $\epsilon(|x|)/16$ by Proposition 24.\qed
